\documentclass[11pt]{article}
\usepackage[margin=1in]{geometry}
\usepackage{amsmath,amssymb,amsthm}
\usepackage{graphicx}
\usepackage{hyperref}
\hypersetup{hidelinks}

\newtheorem{theorem}{Theorem}
\newtheorem{lemma}{Lemma}
\newcommand{\R}{\mathbb R}
\newcommand{\Hent}{\mathcal H}
\newcommand{\Uclass}{\mathcal U}

\title{A Universal and Flux-Sensitive Entropy Lower Bound\\for Weakly Genuinely Nonlinear Scalar Laws}
\author{Run \texttt{fa\_banach\_001}, agent \texttt{agent\_lane\_05}}
\date{August 11, 2026}

\begin{document}
\maketitle

\section*{Source problem}

In arXiv:1912.00219, Capuani, Dutta, and Nguyen consider
\[
 u_t+f(u)_x=0
\]
on \(\R\).  For
\[
 \Uclass_{[L,M]}=\{u_0\in L^\infty(\R):
       \operatorname{supp}u_0\subset[-L,L],\ \|u_0\|_\infty\leq M\},
\]
they obtain an entropy upper bound for \(S_T(\Uclass_{[L,M]})\) when
\(f\in C^2\) is weakly genuinely nonlinear, meaning that \(f\) is affine on
no nonempty open interval.  They ask for sharp \(L^1\)
\(\varepsilon\)-entropy estimates for general fluxes in this class.

\begin{center}
\includegraphics[width=0.95\linewidth]{figures/open_question_crop.png}
\end{center}

The result below supplies a general lower estimate and a sharper
flux-sensitive version.  It does not identify the sharp rate for every such
flux, so its status is partial.

\section*{Statement}

Write \(\Hent_\varepsilon(K\mid L^1)\) for the base-two logarithm of the
least number of \(L^1\)-balls of radius \(\varepsilon\) covering \(K\).

\begin{theorem}[Universal lower bound]\label{thm:universal}
Let \(f\in C^2(\R)\) be weakly genuinely nonlinear.  For every
\(L,M,T>0\), there are constants \(c>0\) and \(\varepsilon_0>0\), depending
on \(f,L,M,T\), such that
\[
 \Hent_\varepsilon\bigl(S_T(\Uclass_{[L,M]})\mid L^1(\R)\bigr)
 \geq \frac{c}{\varepsilon}
 \qquad(0<\varepsilon<\varepsilon_0).
\]
\end{theorem}

The proof gives the following scale-sensitive strengthening.  After replacing
\(f(u)\) by \(f(u)-f'(0)u\), which only translates solutions, suppose that
for a given small \(\varepsilon\) there are \(a=a_\varepsilon\),
\(\delta=A\varepsilon/L\), and \(K=K_\varepsilon>0\) such that
\begin{enumerate}
\item \([a,a+\delta]\subset(-M,M)\), and \(f''\) has one strict sign there;
\item \(\max_{[a,a+\delta]}|f''|\leq K\);
\item the interval joining \(0\) to \([a,a+\delta]\) is contained in a
fixed sufficiently small neighbourhood of zero.
\end{enumerate}
Here \(A>0\) is a universal sufficiently large constant.  Whenever
\(TK\delta\) is small relative to \(L\), the same construction yields
\begin{equation}\label{eq:window}
 \Hent_\varepsilon\bigl(S_T(\Uclass_{[L,M]})\mid L^1(\R)\bigr)
 \geq c_0\frac{L^2}{TK\varepsilon},
\end{equation}
with a universal numerical \(c_0>0\).  Thus locally flatter windows carry
more independent pulses.  For example, if windows of length comparable to
\(\varepsilon\) satisfy \(K_\varepsilon\lesssim
\varepsilon^{m-1}\), then \eqref{eq:window} gives an
\(\varepsilon^{-m}\) lower bound.

\section*{Independent-pulse packing lemma}

\begin{lemma}\label{lem:pulses}
Assume \(f'(0)=0\).  Let \([a,a+\delta]\subset(-M,M)\) be an interval on
which \(f''\) has one strict sign and \(|f''|\leq K\).  Assume also that all
Riemann-wave speeds between the states \(0\) and \([a,a+\delta]\) have
absolute value at most \(L/(16T)\).  If \(\delta\) and
\(q:=8TK\delta\) are sufficiently small compared with \(L\), then
\(S_T(\Uclass_{[L,M]})\) contains a family \(\{v_\sigma\}_{\sigma\in\mathcal C}\)
such that
\[
 |\mathcal C|\geq 2^{N/8},\qquad
 N\geq \frac{L}{16q},\qquad
 \|v_\sigma-v_\tau\|_1\geq \frac{L\delta}{64}
 \quad(\sigma\ne\tau).
\]
\end{lemma}

\begin{proof}
Place a plateau of height \(a\) on \([-3L/4,3L/4]\), with value zero
outside.  Inside \([-L/4,L/4]\), divide the line into \(N\) blocks of
length \(4q\).  In each block choose a fixed subinterval \(I_j\) of length
\(q\).  For \(\sigma=(\sigma_1,\ldots,\sigma_N)\in\{0,1\}^N\), define
\[
 u_\sigma(0,x)=a+\delta\sum_{j=1}^N\sigma_j\mathbf1_{I_j}(x)
\]
on the plateau and set it equal to zero outside the plateau.  These initial
data belong to \(\Uclass_{[L,M]}\).

Every wave created by a pulse takes values in \([a,a+\delta]\).  Its speed
differs from \(f'(a)\) by at most \(K\delta\), by the mean-value theorem
(the same bound applies to shock speeds).  During time \(T\), therefore,
each edge moves relative to \(f'(a)T\) by at most
\(TK\delta=q/8\).  The pulse wave fans do not meet one another because the
blocks have length \(4q\).  The waves generated at the two ends of the large
plateau stay away from the pulse zone: its initial buffer is \(L/2\), while
the assumed speed bound permits closing by at most \(L/8\).

Consequently, up to time \(T\), the solution in the cone of the \(j\)-th
block depends only on \(\sigma_j\).  If \(\sigma_j=1\), comparison with the
constant solution \(a\), together with conservation of the compactly
supported mass difference, gives
\[
 \int_{\text{the }j\text{-th cone}}|S_Tu_\sigma-a|\,dx=\delta q.
\]
If \(\sigma_j=0\), that contribution is zero.  Thus two codewords differing
in \(d_H(\sigma,\tau)\) coordinates satisfy
\begin{equation}\label{eq:hamming}
 \|S_Tu_\sigma-S_Tu_\tau\|_1
 \geq d_H(\sigma,\tau)\delta q.
\end{equation}

The elementary greedy binary-code bound supplies
\(\mathcal C\subset\{0,1\}^N\) with
\(|\mathcal C|\geq2^{N/8}\) and minimum Hamming distance at least \(N/4\)
for all sufficiently large \(N\).  Since the blocks fit in an interval of
length \(L/2\), one may take \(N=\lfloor L/(8q)\rfloor\), and for small
\(q\) this is at least \(L/(16q)\).  Combining this with
\eqref{eq:hamming} proves the stated separation.
\end{proof}

\section*{Proof of the results}

Subtracting the affine term \(f'(0)u\) from the flux corresponds to a common
spatial translation at time \(T\), hence preserves all \(L^1\) distances.
We may assume \(f'(0)=0\).

For the window estimate, take \(A\) large enough that
\(L\delta/64>2\varepsilon\).  Lemma~\ref{lem:pulses} then gives a
\(2\varepsilon\)-separated family of cardinality at least \(2^{N/8}\), and
therefore
\[
 \Hent_\varepsilon(S_T(\Uclass_{[L,M]})\mid L^1)
 \geq \frac{N}{8}
 \geq c_0\frac{L}{TK\delta}
 =c_0'\frac{L^2}{TK\varepsilon}.
\]

For Theorem~\ref{thm:universal}, continuity of \(f'\) at zero gives a
neighbourhood \((-\eta,\eta)\) in which all characteristic and chord speeds,
after the normalization, are smaller than \(L/(16T)\).  Weak genuine
nonlinearity says that \(\{u:f''(u)\ne0\}\) is dense.  Choose one
\(a\in(-\eta/2,\eta/2)\) with \(f''(a)\ne0\).  By continuity there is a
fixed interval \([a,a+r]\subset(-\eta,\eta)\cap(-M,M)\) on which \(f''\)
has one strict sign and \(|f''|\leq K\) for a fixed finite \(K\).
For all sufficiently small \(\varepsilon\), the choice
\(\delta=A\varepsilon/L<r\) satisfies the lemma.  Formula
\eqref{eq:window}, with this fixed \(K\), proves the universal
\(c/\varepsilon\) bound.

\section*{Upgrade attempts and obstruction to a full answer}

Five focused routes were checked.

\begin{enumerate}
\item A fixed locally curved state window yields Theorem~\ref{thm:universal}.
\item Allowing the window to depend on \(\varepsilon\) yields the stronger
curvature-sensitive estimate \eqref{eq:window}; it recovers the usual
polynomial exponents and also applies along selected scales for fluxes with
infinitely many inflection points.
\item Adapting the lower construction of Ancona--Glass--Nguyen confirms that
one-sided controllability gives the same mechanism when curvature keeps one
sign from the zero state, but that hypothesis is absent in the target.
\item Trying to match the source's generalized-variation upper bound fails
at a precise point: weak genuine nonlinearity only makes
\(\{f''\ne0\}\) open dense.  Its sign components near a flat accumulation
point may all be much shorter than the required amplitude scale
\(\delta\asymp\varepsilon\).  The independent-pulse lower construction then
cannot access the flatness that controls the global convex-envelope modulus.
\item Allowing pulses to cross many curvature-sign changes does not repair
this.  The Riemann solution is governed by convex/concave envelopes and can
erase the small oscillatory curvature; the disjoint backward-controllability
argument no longer supplies a matching packing.
\end{enumerate}

Thus the packet proves a general lower bound and a useful conditional
refinement, but a sharp two-sided rate for every weakly genuinely nonlinear
\(C^2\) flux remains open.

\section*{Novelty and literature check}

The cheap run indexes were searched for arXiv \texttt{1912.00219} and the
core phrases concerning weak genuine nonlinearity, entropy, generalized
variation, and scalar conservation laws; no duplicate packet or attempt was
found.  Bounded primary-source searches found the source paper and
Ancona--Glass--Nguyen, arXiv:1806.07758, which treats strictly convex fluxes
and a single polynomially degenerate inflection point.  No later arXiv paper
claiming the general sharp estimate, or this local-window extension, was
found.  The result should therefore receive both PDE-proof and novelty review.

\begin{thebibliography}{9}
\bibitem{CDN}
R. Capuani, P. Dutta, and K. T. Nguyen,
\emph{Metric entropy for functions of bounded total generalized variation},
arXiv:1912.00219.

\bibitem{AGN}
F. Ancona, O. Glass, and K. T. Nguyen,
\emph{On Kolmogorov entropy compactness estimates for scalar conservation
laws without uniform convexity}, arXiv:1806.07758; SIAM J. Math. Anal. 51
(2019), 3020--3051.
\end{thebibliography}

\end{document}
